\subsection{Setup}\label{sec:setup}
\paragraph{Datasets.} We use four image datasets of increasing difficulty: MNIST~\citep{lecun1998mnist}, Fashion-MNIST~\citep{xiao2017fmnist} (the main benchmark), SVHN~\citep{netzer2011svhn} (natural photographs of house numbers) and CIFAR-10~\citep{krizhevsky2009cifar} (natural object images). Section~\ref{sec:real} adds realistic benchmarks with a pretrained ResNet-18: PACS, Office-Home and CIFAR-100. Federations are built from them following the non-IID taxonomy used in the survey~\citep{benali2026survey}, which lists image classification with rotations and label swaps as the standard CFL benchmark~\citep[Table~6]{benali2026survey}.

\paragraph{Heterogeneity.} Each client belongs to a latent concept group. \emph{Concept shift on features} is a rotation of the inputs (by $0^\circ,90^\circ,180^\circ,270^\circ$, or by $0,\theta,2\theta,3\theta$ in the strength sweep). \emph{Concept shift on targets} is a label permutation that swaps two pairs of classes per group. On top of this, every client has \emph{target skew} (Dirichlet label proportions, $\alpha=0.3$, or $\alpha=0.1$ in the label-only scenario) and, where indicated, \emph{quantity skew} (log-normal dataset sizes with $\sigma=1$, clipped to $[30,3200]$, mean 400). Table~\ref{tab:scenarios} lists all scenarios. Each client has a local test set drawn with its own label proportions and transformation. We also evaluate every model on a class-balanced test set of its client's concept (\emph{concept accuracy}), which measures generalisation beyond the local label mix.

\begin{tableorg}[!htbp]
\centering
\caption{Scenarios. $N$: clients; groups: number of concept groups; QS: quantity skew; part.: fraction of clients per round.}\label{tab:scenarios}
\small
\resizebox{\textwidth}{!}{%
\begin{tabular}{llllcl}
\toprule
Scenario & Dataset(s) & Concept groups & $N$ & QS & Other \\
\midrule
Rotation & FMNIST, MNIST, SVHN, CIFAR-10 & 4 rotations ($90^\circ$ steps) & 40 & -- & \\
Label swap & FMNIST, MNIST, SVHN, CIFAR-10 & 4 labelings (2 swapped pairs) & 40 & -- & \\
Rotation + QS & FMNIST & 4 rotations & 40 & \checkmark & \\
Mixed + QS & FMNIST & 2 rotations $\times$ 2 labelings & 40 & \checkmark & \\
Label skew only & FMNIST & 1 (true $K=1$) & 40 & \checkmark & $\alpha=0.1$ \\
Minority groups & FMNIST & 4 rotations, 25/9/4/2 clients & 40 & -- & \\
Strength $\theta$ & FMNIST & 4 rotations, $\theta\in\{15^\circ,30^\circ,45^\circ\}$ & 40 & -- & \\
Groups $K$ & FMNIST & $K\in\{2,4,8\}$ labelings & 40 & -- & \\
Partial participation & FMNIST & rotation / label swap & 40 & -- & part.\ $=0.3$ \\
Scale & FMNIST & 4 rotations & 200 & -- & 200 samples/client \\
\bottomrule
\end{tabular}}
\end{tableorg}

\paragraph{Training.} All methods use the same CNN (two convolutional and two fully connected layers; $d=46{,}730$ parameters for $28\times28$ inputs and $65{,}962$ for $32\times32$ colour images), SGD with learning rate $0.05$, batch size 32, one local epoch per round and 50 communication rounds. Methods with a warm-up (DisCo-CFL, FL+HC, FeSEM, Oracle) spend 10 of the 50 rounds on FedAvg. \emph{Local} trains each client alone for 50 epochs. The six Fashion-MNIST scenarios of Table~\ref{tab:main} and the MNIST, SVHN and CIFAR-10 scenarios are run with five seeds, the sweeps and ablations with three. A seed changes the federation, the initialisation and the sketch. Section~\ref{sec:cost} reports wall-clock time and communication for every method.

\paragraph{Baselines.} We compare with twelve methods that cover all three families of the taxonomy, the three closest class-wise or frozen-model methods, and the main alternatives to clustering:
\begin{itemize}[leftmargin=*,itemsep=1pt]
\item \emph{Single model / personalisation:} FedAvg~\citep{mcmahan2017fedavg}; Local; FedAvg-FT (FedAvg followed by one epoch of local fine-tuning, a standard personalisation baseline).
\item \emph{Client-side:} IFCA~\citep{ghosh2020ifca}.
\item \emph{Server-side:} CFL/MTCFL~\citep{sattler2021cfl} (recursive cosine bipartitioning, triggered when the norm of the aggregated update falls below $\varepsilon_1=0.9$ times the mean client-update norm); FL+HC~\citep{briggs2020flhc} (Ward clustering of local updates after the warm-up); FeSEM~\citep{long2023fesem} (EM-like assignment to the nearest of $K$ centres in parameter space, with the same warm-up and a $k$-means++ initialisation on the clients' local models).
\item \emph{Metadata-based:} PACFL~\citep{vahidian2023pacfl} (principal angles between client data subspaces, $p=3$).
\item \emph{Closest methods}, reimplemented from the papers with the same model, optimiser, data and number of rounds as all other methods: StoCFL~\citep{zeng2025stocfl} (normalised full-batch gradients at the initial model, greedy merging while the cosine exceeds $\tau_{\rm S}=0.5$, bi-level training with a proximal weight $0.05$; two local epochs per round, one for the cluster model and one for the global model); StoCFL-W, the same with the gradients taken at the 10-round warm-up model that DisCo-CFL uses, so that the two differ only in the signal and the rule; FedCCFA~\citep{chen2024fedccfa} (shared feature extractor, per-class clustering of balanced classifier rows by DBSCAN on the MADD distance with radius $0.1$ and one sample per core, classifier averaging within the per-class clusters, and prototype alignment; each client predicts with its personal model, i.e.\ the shared extractor and its classifier after the local classifier epoch of the protocol); and FedDAG-C, the clustering stage of FedDAG~\citep{pramanik2026feddag} (gradient similarity after two local epochs fused with frequency-weighted principal angles between class-wise data subspaces, average linkage with the true $K$, then FedAvg within clusters). FedCCFA clusters each class separately; for ARI we use the partition of clients induced by identical cluster memberships over all classes. The deviations from the original papers (a single SGD optimiser for all methods; for FedDAG-C, a fixed fusion weight instead of its adaptive weighting and the true $K$ instead of an index-based choice) are documented in the code.
\item \emph{Oracle:} FedAvg inside the ground-truth groups (an upper reference).
\end{itemize}
IFCA, FL+HC, FeSEM, PACFL and FedDAG-C are \emph{given the true number of groups} $K$, which favours them. DisCo-CFL, MTCFL, StoCFL and FedCCFA find $K$ on their own. Thresholds of the baselines that find $K$ were selected on the same pilot runs as those of DisCo-CFL: the MTCFL thresholds and the FeSEM initialisation in the original study, i.e.\ before the correction described below and not re-selected afterwards (an asymmetry in favour of DisCo-CFL, whose threshold was re-selected), and, after the correction, the StoCFL threshold ($\tau_{\rm S}\in\{0.1,0.3,0.5,0.7,0.8,0.9\}$, at the initial and at the warm-up model) and the FedCCFA radius ($\{0.02,0.05,0.1,0.2,0.4\}$). No StoCFL threshold reaches a mean pilot ARI above $0.1$, so we keep the published default $0.5$; for FedCCFA the published default $0.1$ has the highest mean pilot ARI. \emph{Budgets.} All methods use 50 rounds of one local epoch per client, except StoCFL (two epochs per round) and FedDAG-C (two warm-up epochs); warm-up rounds count towards the 50. Hyperparameters of the baselines were otherwise taken from their papers, without further tuning for any method.

\paragraph{DisCo-CFL.} Sketch dimension $p=1024$, at most 200 samples per class for the signature, classes with fewer than 4 samples not reported, $\tau=0.9$, $\lambda=0.5$, $e_{\min}=0.05$, low-evidence threshold $0.3$, $z=2$, per-sample clipping $C_{\rm clip}=0.5$. These values, and the design choices of bottleneck linkage and evidence-ordered merging, were fixed during pilot runs on seed 0 of the Rotation, Label swap, Mixed + QS and label-only configurations using clustering quality \emph{measured against the ground-truth groups}, which a deployment would not have. They are kept identical for all scenarios, seeds and datasets. No other seed, sweep or dataset was used for design. Per-sample clipping ($C_{\rm clip}$ chosen on seed 0 only, from $\{0.5,1,2\}$ in all scenarios and additionally $\{0.1,0.25\}$ in the label-only scenario) was added after a first sensitivity study showed over-splitting caused by heavy-tailed per-sample gradients. The run without clipping is reported as an ablation. \emph{Correction and re-selection of the threshold.} After submission we found an implementation error in the training primitive shared by all methods: a model vector passed to several clients in the same round was modified in place, so that clients of the same cluster effectively trained one after another instead of in parallel. This affected every method and every result of the original manuscript. All experiments reported here were rerun with the corrected code. With the corrected code the threshold $\tau=0.8$ merged groups on the pilot runs (ARI $0.60$ on Rotation and $0.64$ on Mixed + QS), and we re-selected it on the same pilot runs (seed 0 of the four configurations, $\tau\in\{0.70,0.75,0.80,0.85,0.90\}$, $T_0\in\{10,20,30\}$). $\tau=0.9$ with $T_0=10$ gives ARI $0.93$, $1.00$, $1.00$ and $1.00$ and was adopted; the other constants were not changed. \emph{Development and evaluation runs.} The pilot runs (seed 0 of Rotation, Label swap, Mixed + QS and label-only, and seed 0 of PACS for the learning rate of Section~\ref{sec:real}) are development runs. We report every result both on all runs and on the held-out runs only, which exclude them. On the colour datasets we also report DisCo-T30, which uses 30 warm-up rounds (Section~\ref{sec:datasets}). Section~\ref{sec:sens} reports the sensitivity to $\tau$ and $\lambda$.

\paragraph{Metrics.} \emph{Performance:} mean local test accuracy (Acc), with the prior correction of Section~\ref{sec:pc} applied identically to every method (for methods that predict with per-client models, such as Local, FedAvg-FT and FedCCFA, the cluster prior equals the client's own prior and the correction has no effect), and concept accuracy. \emph{Fairness:} accuracy of the worst 10\% of clients (Acc@10\%), accuracy of the worst concept group (Worst-grp), Jain's index. \emph{Clustering:} ARI with respect to the groups, number of clusters $K$, and the \emph{conflict-merge rate}: the fraction of conflicting client pairs (Definition~\ref{def:concept}) that end up in the same cluster. ARI penalises the legitimate placement of compatible clients (e.g.\ a client of a label-swap group that holds none of the swapped classes), whereas the conflict-merge rate measures exactly the harmful errors. \emph{Statistics.} Runs with different seeds of the same scenario share the scenario, so (scenario, seed) pairs are not independent. We report, per dataset, means with 95\% $t$-intervals over seeds and \emph{failure rates} (fraction of runs with ARI below $0.5$ when there are several groups). For comparisons across scenarios we use (i) a hierarchical bootstrap that resamples scenarios and then seeds within scenarios~\citep{saravanan2020hierarchical} ($10^4$ replicates) for the mean difference between DisCo-CFL and each baseline, and (ii) one-sided Wilcoxon signed-rank tests on scenario-level means, one value per scenario~\citep{demsar2006statistical}. Both are reported on all runs and on held-out runs.
